\section{Introduction}
\label{sec:introduction}

An economy can have a high proportion of informal workers without producing an equally high proportion of its output informally. The distinction is particularly consequential for Colombia. The national household-survey series used in this paper places labor informality at approximately 54\% of employment in 2026. Yet the relative productivity of work performed in small, unregistered, or own-account activities differs from that of registered wage employment. Translating a headcount into a share of value added requires information about those differences and a definition of what constitutes informal production.

The available indicators measure different objects. Labor informality counts \emph{workers}; DANE's multidimensional business-informality indicator counts \emph{establishments} and reported 90.2\% of microbusinesses as informal in 2024 \citep{DANE_IMIE}. The statistical office's estimate of the non-observed economy, 22.4\% of gross production in its 2015-base compilation, encompasses informal, underground and illegal activities as well as statistical adjustments \citep{DANE_SIECI}. The present study instead estimates the share of \emph{measured value added associated with informal work} under an explicit two-sector production model. These figures should not be equated: their units, denominators and coverage are different.

Dynamic general equilibrium (DGE) models are widely used to infer the size of the informal or shadow economy. Calibrated approaches fix its level at an external benchmark and model the subsequent dynamics \citep{ElginOztunali2012,SolisGarciaXie2017}. Bayesian DSGE approaches estimate latent sectoral activity jointly with structural shocks \citep{OrsiRaggiTurino2014}. Both require an identification argument, because aggregate observations alone do not always distinguish a large informal sector from a small one: different shares can be accommodated by changes in unobserved sectoral productivity and in the implied steady state. This paper supplies that argument with a direct measurement.

We investigate this problem with a deliberately parsimonious two-sector model. Competitive formal firms employ capital and labor; informal activities are represented by a labor-only technology, following a familiar strand of the literature \citep{IhrigMoe2004,ElginOztunali2012}. A household allocates work between the two activities, accumulates capital and chooses consumption. Aggregate shocks, a formality wedge and informal relative productivity generate business-cycle fluctuations. We estimate the model with a Kalman-filter likelihood and Bayesian Markov chain Monte Carlo methods. The macroeconomic data come from DANE's 2015-reference national accounts, with reconstructed quarterly observations before 2005. Total employment and labor informality are drawn from the GEIH-universe ICML series, which reconstructs the national labor market from 1993 \citep{VeraJaramillo2025ICML,VeraJaramillo2026Zenodo}.

The additional identifying information is a quarterly series of \emph{relative labor earnings}, constructed from DANE household-survey microdata since 2007. The ratio distinguishes worker groups and is adjusted for statutory non-wage employment costs on the formal side. Under the baseline production and measurement assumptions, the informal-to-formal output ratio equals the ratio of workers multiplied by relative value added per worker. The latter can be related to the survey earnings ratio and the formal capital share. This yields a transparent link between data and the quantity of interest (Proposition~\ref{prop:share}), in which every assumption is explicit and can be varied.

The paper makes four contributions. First, it provides a quarterly estimate of informal output for Colombia from 1993 to 2026, built on all the observed and reconstructed information available, together with the other quantities the model delivers: informal and formal GDP, relative productivity and the structural shocks behind them. Second, it shows through likelihood profiles, prior--posterior comparisons and simulations with known parameters that aggregate data alone do not determine the size of the informal sector, which explains why calibrated and estimated DGE models in the literature can reach very different levels. Third, it shows that relative earnings solve this problem for the object of interest: with them, the estimated \emph{path} of informal output is invariant to the model's steady-state share, and simulations recover the true share with high precision. Fourth, the data, code and diagnostics are public, so the series can be updated each quarter and re-estimated under alternative definitions, technologies and wage-cost conversions.

Under the preferred specification, informal work accounts for approximately 30.3\% of measured output in 1993 and 20.0\% in 2026Q2. The reduction parallels the decline in the labor-informality rate from approximately 69\% to 54\%; at the end of the sample, the implied informal value added per worker is about one-fifth of its formal counterpart. The estimate falls within the range of published Colombian shadow-economy estimates, close to the enhanced MIMIC series in 2009--2020 and to the Medina--Schneider series in 2008--2015 \citep{SchneiderBuehnMontenegro2010,MedinaSchneider2019,AsllaniDellAnnoSchneider2024,ElginKoseOhnsorgeYu2021}, although those series measure related but not identical concepts.

The paper is a measurement exercise: it reports the level of informal output together with the evidence and the accounting assumptions that determine it, so that each can be checked and improved.

The remainder of the paper reviews measurement approaches (Section~\ref{sec:background}), describes the data and their provenance (Section~\ref{sec:data}), presents the model and estimation (Sections~\ref{sec:model}--\ref{sec:estimation}), examines identification (Section~\ref{sec:identification}), and reports the estimates (Section~\ref{sec:results}). Section~\ref{sec:discussion} discusses their interpretation and Section~\ref{sec:limitations} the extensions that would sharpen them; Section~\ref{sec:conclusions} concludes.
